
%->8-------------------------------------------------------------------------------------------8<-%
\usepackage[brazil]{babel}
\usepackage[utf8]{inputenc}
\usepackage[T1]{fontenc}
%->8-------------------------------------------------------------------------------------------8<-%

%->8-------------------------------------------------------------------------------------------8<-%
\usetheme{Darmstadt}
%->8-------------------------------------------------------------------------------------------8<-%

%->8-------------------------------------------------------------------------------------------8<-%
\usepackage{graphicx}
\usepackage{float}
%->8-------------------------------------------------------------------------------------------8<-%

%->8-------------------------------------------------------------------------------------------8<-%
\usepackage{geometry}
\geometry{left=0.4cm, right=0.4cm}
\usepackage{csquotes}
\usepackage{upquote}
\usepackage{listings}
\usepackage{hyperref}
\setbeamertemplate{navigation symbols}{}
\setbeamertemplate{footline}{\quad\hfill\normalsize{\insertframenumber}\strut\quad\,}
%->8-------------------------------------------------------------------------------------------8<-%

%->8-------------------------------------------------------------------------------------------8<-%
\usepackage[style=numeric, backend=biber]{biblatex}
\addbibresource{outer-crust/repertoire.bib}
\graphicspath{{figures}}
%->8-------------------------------------------------------------------------------------------8<-%

%->8-------------------------------------------------------------------------------------------8<-%
\lstnewenvironment{code}[1][]{
  \lstset{
    basicstyle=\ttfamily,
    columns=flexible,
    breaklines=true,
    breakatwhitespace=true,
    frame=none,
    basewidth=0.5em,
    aboveskip=13pt,
    belowskip=0pt,
    #1
  }
}{}

\newcommand{\excerptpage}[2]{
  \begin{center}
    \textit{#1}\\
        --- #2.
  \end{center}
}

\newcommand{\bigtitlepage}[1]{%
  \section{#1}
  \begin{center}
    \vspace*{\fill}
    {\Huge #1}
    \vspace*{\fill}
  \end{center}
}
%->8-------------------------------------------------------------------------------------------8<-%

% Math
%->8-------------------------------------------------------------------------------------------8<-%
\usepackage{amsmath, amssymb, calculation, mathdots, scalerel, mathtools}
\usepackage{../calculation}

  % Matrices and Vectors Related Commands
  %->8-------------------------------------------------------------8<-%
  \newenvironment{jmatrix}{\left[\begin{matrix}}{\end{matrix}\right]}
  \newcommand{\norm}[1]{\left\lVert#1\right\rVert}
  \newcommand{\fnorm}[1]{\left\lVert#1\right\rVert_\mathsf{F}}
  \newcommand{\dotprod}[2]{\left\langle #1 \, \middle| \, #2 \right\rangle}
  \newcommand{\proj}[2]{\mathrm{proj}_{#1}(#2)}
  \renewcommand{\det}[1]{\mathrm{det}\left(#1\right)}
  \newcommand{\posto}[1]{\mathrm{posto}\left(#1\right)}
  \newcommand{\dist}[1]{\mathrm{dist}(#1)}
  \newcommand{\N}[1]{\mathrm{N}\left(#1\right)}
  \newcommand{\spn}[1]{\mathrm{span}\left(#1\right)}
  \renewcommand{\t}{^{\scriptscriptstyle\mathsf{T}\mkern-1mu}}
  \renewcommand{\dim}[2]{(#1 \times #2)}
  \newcommand{\bigvert}{\rule[-1ex]{0.5pt}{6.8ex}}
  \renewcommand{\vert}{\rule[-1ex]{0.5pt}{2.5ex}}
  \newcommand{\bighori}{\rule[.5ex]{6.8ex}{0.5pt}}
  \newcommand{\hori}{\rule[.5ex]{2.5ex}{0.5pt}}
  %->8-------------------------------------------------------------8<-%

  % Optimization Related Commands
  %->8-------------------------------------------------------------8<-%
  \newcommand{\mintext}{\text{Min}}
  \renewcommand{\min}[1]{\underset{#1}{\mathrm{Min}}\,\,}
  \newcommand{\maxtext}{\text{Max}}
  \renewcommand{\max}[1]{\underset{#1}{\mathrm{Max}}\,\,}
  \newcommand{\argmintext}{\text{Argmin}}
  \newcommand{\argmin}[1]{\underset{#1}{\mathrm{Argmin}}\,\,}
  \newcommand{\argmaxtext}{\text{Argmax}}
  \newcommand{\argmax}[1]{\underset{#1}{\mathrm{Argmax}}\,\,}
  %->8-------------------------------------------------------------8<-%

  % Functions Related Commands
  %->8-------------------------------------------------------------8<-%
  \newcommand{\grad}[1]{\underset{#1}{\nabla}\,\,}
  \newcommand{\hess}[1]{\underset{#1}{\mathcal{H}}\,\,}
  \newcommand{\sderiv}[2]{\dfrac{d #2}{d #1}}
  \newcommand{\bderiv}[2]{\dfrac{d}{d #1}\left(#2\right)}
  \newcommand{\spderiv}[2]{\dfrac{\partial #2}{\partial #1}}
  \newcommand{\bpderiv}[2]{\dfrac{\partial}{\partial #1}\left(\displaystyle#2\right)}
  \newcommand{\dpderiv}[2]{\dfrac{\partial^2\left(\displaystyle#2\right)}{\partial #1^2}}
  \newcommand{\vero}[1]{\mathcal{L}(#1)}
  \newcommand{\cvero}[2]{\mathcal{L}(#1 \mid #2)}
  \newcommand{\prob}[2]{\mathrm{P}_{#1}\left[#2\right]}
  \newcommand{\cprob}[3]{\mathrm{P}_{#1}\left(#2 \mid #3\right)}
  \newcommand{\pfunc}[2]{\mathrm{P}_{#1}\left(#2\right)}
  %->8-------------------------------------------------------------8<-%

  % Number Related Commands
  %->8-------------------------------------------------------------8<-%
  \newcommand{\mul}{\cdot}
  \renewcommand{\deg}{\ensuremath{^\circ}}
  \renewcommand{\sin}[1]{\mathrm{sen}\left(#1\right)}
  \newcommand{\logtext}{\mathrm{log}}
  \renewcommand{\log}[1]{\logtext\left(#1\right)}
  \newcommand{\lntext}{\mathrm{ln}}
  \renewcommand{\ln}[1]{\lntext\left(#1\right)}
  \renewcommand{\exp}[1]{\mathrm{e}^{#1}}
  \newcommand{\abs}[1]{\left\lvert#1\right\rvert}
  %->8-------------------------------------------------------------8<-%

  % Modeling Related Commands
  %->8-------------------------------------------------------------8<-%
  \newcommand{\blackbox}{\raisebox{0.2ex}{\rule{1.2ex}{1.2ex}}}
  \newcommand{\whitebox}{\raisebox{0.2ex}{\fboxsep=0pt\fbox{\rule{1.2ex}{0pt}\rule{0pt}{1.2ex}}}}
  \newcommand{\fim}[2]{\vspace{0.075cm} \dfrac{#1}{#2}}
  %->8-------------------------------------------------------------8<-%

  % Other Math Related Commands
  %->8-------------------------------------------------------------8<-%
  \newcommand{\rep}[1]{\mathrm{rep}\left(#1\right)}
  % \newcommand{\union}{\cup}
  \newcommand{\bigunion}{\bigcup}
  % \newcommand{\intersec}{\cap}
  \newcommand{\bigintersec}{\bigcap}
  % \renewcommand{\gets}{\leftarrow}
  \newcommand{\E}{\Sigma}
  \newcommand{\st}{\,\,|\,\,}
  \newcommand{\bigst}{\,\,\,\,\rule[-0.5cm]{0.5pt}{6.8ex}\,\,\,\,}
  \newcommand{\impliesdown}{\displaystyle\Downarrow}
  % \DeclareMathOperator*{\com}{\scalerel*{,}{\sum}}
  %->8-------------------------------------------------------------8<-%
%->8-------------------------------------------------------------------------------------------8<-%

%->8-------------------------------------------------------------------------------------------8<-%
\title{Raciocínio Equacional para a Derivação de \\ Algoritmos de Álgebra Linear sob a \\ Perspectiva de Otimização}
\author{%
  \textbf{João Victor Lopez Pereira} \\
  \vspace{0.2cm}
  Daniel Sadoc Menasche \\
  \vspace{0.2cm}
  João Antonio Recio Paixão
}
\institute{Instituto de Computação -- UFRJ}
\date{Outubro de 2026}
%->8-------------------------------------------------------------------------------------------8<-%
